%% Datasets documentation — Internship project
%% Continuation of the BOAMP API document (section 1) written by the teammate.
%% Covers: DECP (section 2), processed datasets (section 3),
%% taxonomy (section 4), and renewal-linking output (section 5).

\documentclass[12pt,a4paper]{article}
\usepackage[utf8]{inputenc}
\usepackage[T1]{fontenc}
\usepackage[french]{babel}
\usepackage{geometry}
\usepackage{listings}
\usepackage{xcolor}
\usepackage{booktabs}
\usepackage{array}
\usepackage{hyperref}
\usepackage{parskip}

\geometry{margin=2.5cm}

\lstset{
  basicstyle=\small\ttfamily,
  breaklines=true,
  frame=single,
  backgroundcolor=\color{gray!8},
  keywordstyle=\color{blue},
  stringstyle=\color{red!70!black},
  showstringspaces=false,
}

\begin{document}

% ============================================================
\section{DECP -- Données Essentielles de la Commande Publique}
% ============================================================

\subsection{Présentation générale}

Le jeu de données DECP constitue la source complémentaire principale du projet.
Introduit par le décret du 25 mars 2016, il impose à toute autorité publique
attribuant un marché supérieur à 25\,000~€ de publier un enregistrement
structuré sur \texttt{data.gouv.fr} dans les deux mois suivant la signature.
Ces données sont consolidées par l'AIFE (Agence pour l'Informatique Financière
de l'État) et mises à disposition sous la forme d'un fichier Parquet unique.

Contrairement au BOAMP, qui enregistre des \emph{annonces}, le DECP enregistre
le \emph{contrat signé}. C'est donc la source de référence pour les champs
nécessaires à l'analyse de survie (montant réel, durée déclarée, date de
notification).

\medskip
\textbf{Fichier utilisé}

La version exploitée n'est pas les fichiers JSON annuels officiels (dont la
structure est inconsistante selon les millésimes) mais le fichier Parquet
consolidé publié par le projet communautaire \textit{decp-processing} sur
\texttt{data.gouv.fr} (identifiant de ressource \texttt{11cea8e8-df3e-4ed1-932b-781e2635e432}).
Ce fichier est enrichi par des jointures SIRENE, ce qui ajoute des colonnes
géographiques (commune, département, région) absentes du DECP brut.

\medskip
\textbf{Volume et couverture}

\begin{itemize}
  \item[---] Nombre total d'enregistrements : environ \textbf{3,1 millions} de
    lignes pour l'ensemble de la France et de tous les secteurs
  \item[---] Granularité : un enregistrement correspond à un contrat (ou une
    modification de contrat)
  \item[---] Couverture temporelle dense : à partir de \textbf{2019} (la
    publication était obligatoire depuis 2018 mais la couverture réelle est
    quasi nulle avant 2019)
  \item[---] Corpus retenu après filtrage : \textbf{3\,039 contrats} portant
    sur les marchés numériques en Pays de la Loire (2018--2024), uniquement
    les versions courantes (\texttt{donneesActuelles = True})
\end{itemize}

\medskip
\textbf{Filtres appliqués pour constituer le corpus de travail}

\begin{itemize}
  \item[---] Département de l'acheteur : 44, 49, 53, 72, 85 (Pays de la Loire)
  \item[---] Code CPV : premiers chiffres dans \{48, 72, 32, 35\}
    (marchés numériques / informatiques)
  \item[---] Année de notification : 2015 à 2024
  \item[---] Version courante uniquement : \texttt{donneesActuelles = True}
    (élimine les modifications historiques en doublon)
\end{itemize}

\subsection{Principales données récupérables}

Les informations accessibles sont regroupées en plusieurs catégories.

\medskip
\textbf{Identification du contrat}

\begin{itemize}
  \item[---] \texttt{uid} : identifiant unique global (concaténation de
    \texttt{acheteur\_id} et \texttt{id})
  \item[---] \texttt{id} : identifiant attribué par l'acheteur (non unique
    globalement)
  \item[---] \texttt{nature} : type de contrat (Marché, Accord-cadre, Marché
    subséquent\ldots)
  \item[---] \texttt{donneesActuelles} : booléen indiquant la version courante
    du contrat
  \item[---] \texttt{sourceDataset} : plateforme d'origine de la publication
    (AIFE, PES, marchés-publics.info, e-marchespublics.com\ldots)
\end{itemize}

\medskip
\textbf{Caractéristiques du marché}

\begin{itemize}
  \item[---] \texttt{objet} : intitulé du marché (texte libre, plus concis que
    celui du BOAMP)
  \item[---] \texttt{codeCPV} : code CPV principal ; format structuré, toujours
    présent
  \item[---] \texttt{procedure} : libellé de la procédure standardisé
    (Procédure adaptée, Appel d'offres ouvert, Marché négocié\ldots)
  \item[---] \texttt{montant} : montant du contrat en euros ; champ obligatoire
    du DECP (couverture 98\,\% avant nettoyage)
  \item[---] \texttt{offresRecues} : nombre d'offres reçues ; champ facultatif,
    présent dans 26\,\% des enregistrements du corpus
\end{itemize}

\medskip
\textbf{Informations temporelles}

\begin{itemize}
  \item[---] \texttt{dateNotification} : date de notification du contrat signé ;
    champ obligatoire -- c'est l'\textbf{origine temporelle} de l'analyse de
    survie
  \item[---] \texttt{datePublicationDonnees} : date de publication en
    open data (peut être postérieure de plusieurs mois à la signature)
  \item[---] \texttt{dureeMois} : durée déclarée du contrat en mois ; champ
    obligatoire (couverture 99\,\% dans le corpus)
\end{itemize}

\medskip
\textbf{Informations géographiques et sectorielles}

\begin{itemize}
  \item[---] \texttt{lieuExecution\_code} et \texttt{lieuExecution\_typeCode} :
    code et type du lieu d'exécution (code postal, commune, département ou
    région)
  \item[---] \texttt{acheteur\_departement\_code} : code du département de
    l'acheteur (enrichissement SIRENE) -- utilisé pour le filtre géographique
  \item[---] \texttt{acheteur\_region\_code} / \texttt{acheteur\_region\_nom} :
    région de l'acheteur (enrichissement SIRENE)
\end{itemize}

\medskip
\textbf{Acteurs}

\begin{itemize}
  \item[---] \texttt{acheteur\_id} : SIRET de l'acheteur (14 chiffres) ;
    \textbf{obligatoire} dans le DECP, couverture 100\,\%
  \item[---] \texttt{acheteur\_nom} : nom de l'acheteur (enrichissement SIRENE)
  \item[---] \texttt{acheteur\_categorie} : catégorie de l'acheteur (Commune,
    Groupement de communes, Département, Établissement hospitalier, EPIC\ldots)
  \item[---] \texttt{titulaire\_id} : identifiant du titulaire (SIRET dans la
    plupart des cas)
  \item[---] \texttt{titulaire\_typeIdentifiant} : type d'identifiant du
    titulaire (SIRET, TVA, HORS-UE\ldots)
  \item[---] \texttt{titulaire\_nom} : nom du titulaire (enrichissement SIRENE)
\end{itemize}

\subsection{Comparaison de couverture avec le BOAMP}

\begin{center}
\begin{tabular}{lccc}
\toprule
\textbf{Champ} & \textbf{BOAMP} & \textbf{DECP} &
  \textbf{Source recommandée} \\
\midrule
SIRET acheteur    & 8\,\%   & 100\,\% & DECP \\
Montant (EUR)     & 43\,\%  & 93\,\%  & DECP \\
Durée (mois)      & 34\,\%  & 99\,\%  & DECP \\
Date notification & ---     & 100\,\% & DECP \\
Texte du marché   & 100\,\% & 100\,\% & BOAMP (plus riche) \\
Type d'avis       & 100\,\% & ---     & BOAMP \\
Liens de renouvellement & 46\,\% & ---  & BOAMP (\texttt{annonce\_lie}) \\
\bottomrule
\end{tabular}
\end{center}

\subsection{Exemple d'enregistrement DECP}

\begin{lstlisting}[language={}]
{
  "uid"                    : "214401622000112024K-P0rgdsDS00",
  "id"                     : "2024K-P0rgdsDS00",
  "nature"                 : "Marché",
  "acheteur_id"            : "21440162200011",
  "acheteur_nom"           : "COMMUNE DE SAINT-HERBLAIN (MAIRIE)",
  "acheteur_departement_code" : "44",
  "acheteur_categorie"     : "Commune",
  "titulaire_id"           : "44905091600021",
  "titulaire_nom"          : "AXIANS IDF",
  "objet"                  : "Infogérance du système d'information",
  "codeCPV"                : "72000000",
  "procedure"              : "Procédure adaptée",
  "montant"                : 140000.0,
  "dureeMois"              : 48.0,
  "dateNotification"       : "2024-12-31",
  "datePublicationDonnees" : "2025-01-10",
  "offresRecues"           : null,
  "donneesActuelles"       : true,
  "sourceDataset"          : "aife_dume"
}
\end{lstlisting}

\medskip
\textbf{Cas d'usage dans le projet}

\begin{itemize}
  \item[---] Source principale pour l'analyse de survie : date de notification
    (origine), durée déclarée et montant
  \item[---] Clé SIRET de l'acheteur pour le regroupement et le suivi
    longitudinal des contrats d'un même acheteur
  \item[---] Base pour la construction de la variable cible de renouvellement
    via l'algorithme de liaison (\texttt{decp\_renewal\_links.csv})
  \item[---] Enrichissement de la couverture géographique et sectorielle du
    corpus BOAMP (champs absents ou incomplets dans BOAMP legacy)
\end{itemize}

% ============================================================
\section{Jeux de données traités}
% ============================================================

Les données brutes issues du BOAMP et du DECP sont transformées par un pipeline
de nettoyage (tâche 7) qui produit deux fichiers d'analyse prêts à l'emploi.

% ------------------------------------------------------------
\subsection{\texttt{boamp\_clean.csv} -- Échantillon BOAMP nettoyé}
% ------------------------------------------------------------

\textbf{Contenu}

Échantillon de \textbf{500 avis} BOAMP couvrant la période 2015--2024 en Pays
de la Loire, domaine numérique (CPV 48 / 72 / 32 / 35), équilibrés à 50 avis
par année et mixant les types d'avis (appels d'offres, attributions, etc.).
Le fichier est le résultat de la mise à plat (\textit{flatten}) du champ
\texttt{donnees} et des transformations de nettoyage.

\medskip
\textbf{Volume et structure}

\begin{itemize}
  \item[---] Dimensions : 500 lignes × 44 colonnes
  \item[---] Format : CSV (séparateur virgule, encodage UTF-8)
  \item[---] Couverture temporelle : 2015--2024 (50 avis par année)
\end{itemize}

\medskip
\textbf{Répartition des avis}

\begin{itemize}
  \item[---] \texttt{APPEL\_OFFRE} : 225 avis (45\,\%)
  \item[---] \texttt{ATTRIBUTION} : 217 avis (43\,\%)
  \item[---] \texttt{RECTIFICATIF} : 53 avis (11\,\%)
  \item[---] \texttt{MODIFICATION} / \texttt{PRE-INFORMATION} : 5 avis (1\,\%)
\end{itemize}

\medskip
\textbf{Format des données dans \texttt{donnees}}

\begin{itemize}
  \item[---] \texttt{donnees\_format = legacy} : 460 avis (format BOAMP
    classique, clés françaises : \texttt{IDENTITE}, \texttt{OBJET},
    \texttt{PROCEDURE}\ldots)
  \item[---] \texttt{donnees\_format = eforms} : 40 avis (format UBL européen
    2024, clé racine \texttt{EFORMS}, structure XML imbriquée)
\end{itemize}

\medskip
\textbf{Colonnes ajoutées lors du nettoyage}

\begin{itemize}
  \item[---] \texttt{cpv\_clean}, \texttt{cpv\_div2}, \texttt{cpv\_class4} :
    code CPV normalisé, division (2 chiffres) et classe (4 chiffres)
  \item[---] \texttt{buyer\_siret} : SIRET de l'acheteur extrait de
    \texttt{donnees} (disponible uniquement pour les avis eForms 2024 ; taux
    de renseignement : 8\,\%)
  \item[---] \texttt{buyer\_key} : clé canonique de l'acheteur -- SIRET si
    disponible, sinon nom normalisé (minuscules, sans accents, sans suffixes
    juridiques)
  \item[---] \texttt{amount\_raw}, \texttt{flag\_amount\_zero},
    \texttt{flag\_amount\_tiny}, \texttt{flag\_amount\_ceiling},
    \texttt{amount\_clean} : montant brut, indicateurs de valeurs aberrantes
    et montant nettoyé (taux de renseignement après nettoyage : 43\,\%)
  \item[---] \texttt{duration\_raw\_boamp}, \texttt{flag\_duration\_suspect},
    \texttt{duration\_clean} : durée brute, indicateur et durée nettoyée
    (taux de renseignement après nettoyage : 34\,\%)
  \item[---] \texttt{category\_id}, \texttt{category\_label} : identifiant et
    libellé de la catégorie technologique (cf.\ section~\ref{sec:taxonomy})
\end{itemize}

\medskip
\textbf{Statistiques descriptives des champs quantitatifs}

\begin{center}
\begin{tabular}{lcc}
\toprule
\textbf{Indicateur} & \textbf{Montant (EUR)} &
  \textbf{Durée (mois)} \\
\midrule
Observations valides & 215 & 171 \\
Moyenne              & 814\,356 & 29,4 \\
Médiane              & 430\,070 & 24,0 \\
Min                  & 3\,648   & 1,0 \\
Max                  & 9\,750\,000 & 96,0 \\
\bottomrule
\end{tabular}
\end{center}

% ------------------------------------------------------------
\subsection{\texttt{decp\_clean.csv} -- Corpus DECP nettoyé}
% ------------------------------------------------------------

\textbf{Contenu}

Corpus filtré et nettoyé issu du fichier DECP Parquet, comprenant l'ensemble
des contrats numériques en Pays de la Loire sur la période effective de
couverture dense du DECP.

\medskip
\textbf{Volume et structure}

\begin{itemize}
  \item[---] Dimensions : 3\,039 lignes × 36 colonnes
  \item[---] Format : CSV (séparateur virgule, encodage UTF-8)
  \item[---] Couverture temporelle : 2018--2024 (couverture dense à partir de
    2019 : 361, 433, 506, 574, 561 et 562 contrats par année de 2019 à 2024)
\end{itemize}

\medskip
\textbf{Répartition par nature de contrat}

\begin{itemize}
  \item[---] Marchés : 1\,696 (55,8\,\%)
  \item[---] Accords-cadres : 1\,241 (40,8\,\%)
  \item[---] Marchés subséquents : 99 (3,3\,\%)
  \item[---] Autres : 3 (0,1\,\%)
\end{itemize}

\medskip
\textbf{Répartition géographique}

\begin{itemize}
  \item[---] Loire-Atlantique (44) : 1\,438 contrats (47\,\%)
  \item[---] Maine-et-Loire (49) : 528 contrats (17\,\%)
  \item[---] Sarthe (72) : 508 contrats (17\,\%)
  \item[---] Vendée (85) : 362 contrats (12\,\%)
  \item[---] Mayenne (53) : 203 contrats (7\,\%)
\end{itemize}

\medskip
\textbf{Principales procédures observées}

\begin{itemize}
  \item[---] Procédure adaptée (MAPA) : 1\,201 contrats (39\,\%)
  \item[---] Appel d'offres ouvert : 749 contrats (25\,\%)
  \item[---] Marché négocié sans publicité : 220 contrats (7\,\%)
  \item[---] Marché passé sans publicité ni mise en concurrence : 163 contrats
    (5\,\%)
  \item[---] Autres : 706 contrats (23\,\%)
\end{itemize}

\medskip
\textbf{Colonnes ajoutées lors du nettoyage}

Les mêmes transformations que pour \texttt{boamp\_clean.csv} sont appliquées :

\begin{itemize}
  \item[---] \texttt{buyer\_key} : clé canonique de l'acheteur -- toujours un
    SIRET dans le DECP (100\,\% de couverture sur \texttt{acheteur\_id})
  \item[---] \texttt{cpv\_clean}, \texttt{cpv\_div2}, \texttt{cpv\_class4} :
    code CPV normalisé et agrégats hiérarchiques
  \item[---] \texttt{amount\_raw}, \texttt{flag\_amount\_zero},
    \texttt{flag\_amount\_tiny}, \texttt{flag\_amount\_ceiling},
    \texttt{amount\_clean} : montant nettoyé (taux de renseignement : 93\,\%)
  \item[---] \texttt{duration\_raw\_decp}, \texttt{flag\_duration\_suspect},
    \texttt{duration\_clean} : durée nettoyée (taux de renseignement : 99\,\%)
  \item[---] \texttt{category\_id}, \texttt{category\_label} : catégorie
    technologique (cf.\ section~\ref{sec:taxonomy})
\end{itemize}

\medskip
\textbf{Statistiques descriptives des champs quantitatifs}

\begin{center}
\begin{tabular}{lcc}
\toprule
\textbf{Indicateur} & \textbf{Montant (EUR)} &
  \textbf{Durée (mois)} \\
\midrule
Observations valides & 2\,830 & 3\,014 \\
Moyenne              & 358\,142 & 38,2 \\
Médiane              & 102\,382 & 48,0 \\
1\textsuperscript{er} quartile & 49\,746 & 24,0 \\
3\textsuperscript{e} quartile  & 240\,000 & 48,0 \\
Min                  & 1\,000     & 1,0 \\
Max                  & 9\,750\,000 & 120,0 \\
Mode de la durée     & ---      & 48,0 \\
\bottomrule
\end{tabular}
\end{center}

% ------------------------------------------------------------
\subsection{Règles de nettoyage communes aux deux sources}
% ------------------------------------------------------------

\textbf{Normalisation de la clé acheteur (\texttt{buyer\_key})}

\begin{itemize}
  \item[---] \textbf{Priorité 1 -- SIRET} : si un identifiant numérique de
    14 chiffres est disponible, la clé est \texttt{SIRET:\{14 chiffres\}}
  \item[---] \textbf{Priorité 2 -- Nom normalisé} : minuscules, suppression
    des accents, suppression de la ponctuation, suppression des suffixes
    juridiques (SARL, SA, SAS, Commune de, Mairie de, SDIS, CHU\ldots) ;
    la clé est \texttt{NAME:\{nom normalisé\}}
  \item[---] Table de correspondance inter-sources : \texttt{buyer\_bridge.csv}
    (461 clés canoniques uniques, dont 303 ancrées sur un SIRET)
\end{itemize}

\medskip
\textbf{Nettoyage des montants}

\begin{itemize}
  \item[---] \texttt{flag\_amount\_zero} : valeur = 0 (artefact d'encodage)
  \item[---] \texttt{flag\_amount\_tiny} : $0 < \text{valeur} < 1\,000$~€
    (implausible pour un marché numérique)
  \item[---] \texttt{flag\_amount\_ceiling} : valeur $\geq 10\,000\,000$~€
    (vraisemblablement un plafond d'accord-cadre, non le prix réel)
  \item[---] \texttt{amount\_clean} : NaN si l'un des trois indicateurs est
    vrai ; sinon valeur originale (aucune suppression des données brutes)
\end{itemize}

\medskip
\textbf{Nettoyage des durées}

\begin{itemize}
  \item[---] \texttt{flag\_duration\_suspect} : durée hors de l'intervalle
    $[1, 120]$ mois
  \item[---] \texttt{duration\_clean} : NaN si l'indicateur est vrai, sinon
    valeur originale
  \item[---] Résultats : 0 valeurs suspectes dans DECP ; quelques valeurs
    suspectes dans BOAMP (durées déclarées en jours mal converties)
\end{itemize}

% ============================================================
\section{Taxonomie sectorielle (\texttt{taxonomy.csv})}
\label{sec:taxonomy}
% ============================================================

\textbf{Présentation}

Le fichier \texttt{taxonomy.csv} définit un référentiel de \textbf{10
catégories technologiques} permettant de regrouper les marchés numériques
au-delà des codes CPV bruts (dont il existe environ 9\,000 valeurs distinctes).
Il est utilisé lors de la tâche 7 pour annoter chaque contrat avec un libellé
métier cohérent, aussi bien dans \texttt{boamp\_clean.csv} que dans
\texttt{decp\_clean.csv}.

\medskip
\textbf{Logique d'affectation}

\begin{enumerate}
  \item Correspondance CPV : le préfixe du code CPV normalisé est comparé
    aux préfixes de chaque catégorie ; le préfixe le plus long gagne
    (\textit{longest prefix match})
  \item Repli textuel : si aucun préfixe CPV ne correspond, une recherche de
    mots-clés (insensible à la casse) est effectuée dans le champ
    \texttt{objet}
  \item Valeur par défaut : \texttt{CAT\_UNKNOWN} / \textit{Unknown}
\end{enumerate}

\medskip
\textbf{Référentiel des catégories}

\begin{center}
\begin{tabular}{clcc}
\toprule
\textbf{ID} & \textbf{Libellé} & \textbf{Division CPV} &
  \textbf{Préfixes CPV} \\
\midrule
CAT01 & Software \& Applications          & 48     & 48 \\
CAT02 & IT Services \& Consulting         & 72     & 72 \\
CAT03 & Cybersecurity                     & 35     & 35, 72250, 72310, 72315 \\
CAT04 & Telecom \& Networks               & 32     & 32 \\
CAT05 & Cloud \& Infrastructure           & 48800  & 48800, 72250, 72310 \\
CAT06 & IT Hardware \& Equipment          & 30     & 30, 32420 \\
CAT07 & Digital Workplace \& Collaboration & 48500 & 48500, 72212, 64200 \\
CAT08 & Data \& AI                        & 72300  & 72300, 72322, 72316 \\
CAT09 & GIS \& Mapping                    & 72320  & 72320, 38221 \\
CAT10 & IT Maintenance \& Support         & 50300  & 50300, 72250, 72611 \\
\bottomrule
\end{tabular}
\end{center}

\medskip
\textbf{Distribution des catégories dans le corpus}

\begin{center}
\begin{tabular}{lcc}
\toprule
\textbf{Catégorie} & \textbf{BOAMP (500)} &
  \textbf{DECP (3\,039)} \\
\midrule
IT Services \& Consulting         & 167 (33\,\%) & 1\,321 (43\,\%) \\
Software \& Applications          &  88 (18\,\%) &   608 (20\,\%) \\
Telecom \& Networks               &  80 (16\,\%) &   581 (19\,\%) \\
Cybersecurity                     &  63 (13\,\%) &   337 (11\,\%) \\
Digital Workplace \& Collaboration &  28 (6\,\%) &   110 (4\,\%) \\
Unknown                           &  37 (7\,\%) &     0 (0\,\%) \\
Autres (CAT05--CAT10)             &  37 (7\,\%) &    82 (3\,\%) \\
\bottomrule
\end{tabular}
\end{center}

\medskip
\textbf{Limites connues}

\begin{itemize}
  \item[---] Les codes CPV génériques (ex.\ \texttt{72000000},
    division entière) ne sont pas distingués par le repli textuel si
    l'objet du marché est trop court ou générique ; ils tombent dans
    \texttt{Unknown}
  \item[---] Certaines catégories sont proches (CAT03 Cybersecurity,
    CAT05 Cloud) et partagent des préfixes CPV : la règle du
    \textit{longest prefix} tranche, mais peut induire des erreurs sur les
    marchés transverses
\end{itemize}

% ============================================================
\section{Liens de renouvellement (\texttt{decp\_renewal\_links.csv})}
\label{sec:renewal}
% ============================================================

\textbf{Présentation}

Ce fichier est le \textbf{jeu de données cible} pour l'analyse de survie.
Il contient, pour chaque contrat du corpus DECP nettoyé, soit la durée
observée jusqu'au renouvellement (événement), soit la durée censurée jusqu'à
la fin d'étude (31 décembre 2024).

Ni le BOAMP ni le DECP ne fournissent de champ \texttt{renewal\_of} direct.
Le lien doit donc être reconstruit algorithmiquement.

\medskip
\textbf{Volume et structure}

\begin{itemize}
  \item[---] Dimensions : 3\,039 lignes × 13 colonnes (une ligne par contrat
    du corpus \texttt{decp\_clean.csv})
  \item[---] Format : CSV (séparateur virgule, encodage UTF-8)
\end{itemize}

\medskip
\textbf{Principales données}

\begin{itemize}
  \item[---] \texttt{uid} : identifiant unique du contrat source
    (clé de jointure vers \texttt{decp\_clean.csv})
  \item[---] \texttt{acheteur\_id} : SIRET de l'acheteur
  \item[---] \texttt{cpv\_div2} : division CPV (2 chiffres) utilisée pour le
    regroupement
  \item[---] \texttt{source\_date\_notification} : date de notification du
    contrat source (point zéro de la durée observée)
  \item[---] \texttt{declared\_duration\_months} : durée déclarée
    (fallback : 24 mois si absente ou invalide)
  \item[---] \texttt{renewal\_uid} : identifiant du contrat de renouvellement
    identifié ; \texttt{null} si censuré
  \item[---] \texttt{event} : 1 si un renouvellement a été trouvé, 0 sinon
  \item[---] \texttt{observed\_duration\_months} : durée observée en mois --
    durée jusqu'au renouvellement si \texttt{event=1}, durée jusqu'au
    31/12/2024 si censuré
  \item[---] \texttt{gap\_months} : écart en mois entre la date de notification
    du contrat source et celle du renouvellement
  \item[---] \texttt{text\_similarity} : similarité de Jaccard sur les tokens
    des champs \texttt{objet} des deux contrats
  \item[---] \texttt{match\_score} : score composite de correspondance
    ($0{,}7 \times \text{similarité textuelle} + 0{,}3 \times \text{proximité temporelle}$)
  \item[---] \texttt{is\_censored} : 1 si aucun renouvellement trouvé
    (complément de \texttt{event})
  \item[---] \texttt{censoring\_date} : date de censure (31/12/2024) pour les
    observations censurées
\end{itemize}

\medskip
\textbf{Algorithme de liaison (prototype, semaine 3)}

La liaison est effectuée groupe par groupe : pour chaque paire
(acheteur, division CPV), les contrats sont triés par date, et pour chaque
contrat source on recherche un contrat candidat B vérifiant simultanément :

\begin{enumerate}
  \item $\text{dateNotification}(B) > \text{dateNotification}(\text{source})$
  \item $|\,\text{gap}(B) - \text{dureeMois(source)}\,| \leq 6$ mois
    (fenêtre temporelle de $\pm 6$ mois autour de la durée déclarée)
  \item $\text{Jaccard}(\text{objet source}, \text{objet } B) \geq 0{,}30$
\end{enumerate}

Si plusieurs candidats passent ces filtres, le meilleur est sélectionné par le
score composite décrit ci-dessus.

\medskip
\textbf{Résultats du prototype}

\begin{center}
\begin{tabular}{lc}
\toprule
\textbf{Indicateur} & \textbf{Valeur} \\
\midrule
Contrats analysés               & 3\,039 \\
Renouvellements liés (événements) & 252 (8,3\,\%) \\
Contrats censurés               & 2\,787 (91,7\,\%) \\
Durée médiane observée (tous)   & 30,9 mois \\
Durée médiane observée (événements) & 33,8 mois \\
Similarité textuelle médiane (événements) & 0,60 \\
\bottomrule
\end{tabular}
\end{center}

\medskip
\textbf{Limites et perspectives}

\begin{itemize}
  \item[---] Taux de liaison actuel (8,3\,\%) est un \textbf{prototype de
    référence} ; la cible est 40--60\,\% après amélioration
  \item[---] Piste 1 : affiner le regroupement par classe CPV à 4 chiffres
    plutôt que 2 (plus précis, réduira les faux positifs inter-catégories)
  \item[---] Piste 2 : contrainte un-à-un (chaque contrat source ne peut être
    lié qu'à un seul renouvellement, et vice-versa)
  \item[---] Piste 3 : évaluation manuelle de la précision et du rappel sur un
    échantillon de 50 paires, pour calibrer les seuils de similarité et de
    fenêtre temporelle
  \item[---] Piste 4 : exploiter le champ BOAMP \texttt{annonce\_lie} pour les
    contrats croisables entre les deux sources
\end{itemize}

\end{document}
